\documentclass[14pt]{extarticle}
\usepackage{fontspec}
\usepackage{polyglossia}
\usepackage{amsmath}
\usepackage{amssymb}
\usepackage{xcolor}
\usepackage{fancyhdr}
\usepackage{graphicx}
\usepackage{listings}
\usepackage[most]{tcolorbox}
% \usepackage{setspace}
\usepackage{pifont}
\usepackage{enumitem}
\usepackage{titlesec}
\usepackage[bottom]{footmisc}
\usepackage{titling}
\usepackage{minted}
\usepackage{etoolbox}
\usepackage{array}
\usepackage{longtable} % page-breakable tables (must load before polyglossia/bidi RTL setup)
\usepackage{extsizes}
\usepackage{float}
\usepackage{booktabs}
\usepackage{hyperref}
\usepackage[onehalfspacing]{setspace}
\newif\ifwhiteblack
\ifdefined\setwhiteblack
  \setwhiteblack
\fi

\ifwhiteblack
  \usepackage[a4paper, landscape, top=1.25cm, bottom=1.25cm, left=1.25cm, right=1.25cm, headheight=1.5cm, headsep=0.5cm]{geometry}
\else
  \usepackage[a4paper, top=2.7cm, bottom=2.7cm, left=2.5cm, right=2.5cm, headheight=1.5cm, headsep=0.5cm]{geometry}
\fi
\usepackage{tikz}
\usepackage{datetime2}

\usetikzlibrary{automata, positioning, arrows.meta, arrows}
\providecommand{\LTcaptype}{table}

\newfontfamily\emoji{Segoe UI Emoji}

\pagestyle{fancy}

\setmainlanguage[numerals=western]{arabic}
\setotherlanguage{english}
\newfontfamily\arabicfont[Script=Arabic]{Calibri}
\newfontfamily\arabicfonttt[Script=Arabic]{Courier New}



\ifwhiteblack
\lstset{
  language=[Sharp]C,
  numbers=left,
  stepnumber=1,
  numbersep=8pt,
  frame=single,
  basicstyle=\ttfamily\small,
  keywordstyle=\bfseries,
  stringstyle=\ttfamily,
  commentstyle=\itshape
}
\else
\lstset{
  language=[Sharp]C,
  numbers=left,
  stepnumber=1,
  numbersep=8pt,
  frame=single,
  basicstyle=\ttfamily\small,
  keywordstyle=\color{blue},
  stringstyle=\color{red},
  commentstyle=\color{green!50!black}
}
\fi

\newif\ifdetailed
\ifdefined\setdetailed
  \setdetailed
\fi

\newif\ifwithsols
\ifdefined\setwithsols
  \setwithsols
\fi

\ifwhiteblack
  \AtBeginDocument{
    \LARGE
    \usemintedstyle{bw}
  }
\fi

% Default minted look: compact code that wraps long lines.
% Not applied in black-and-white landscape mode, which already uses \LARGE.
\ifwhiteblack\else
  \setminted{fontsize=\small, breaklines=true}
\fi

% unified tcolorboxes for articles
\ifwhiteblack
\tcbset{colback=white, colframe=black, fonttitle=\bfseries, boxrule=1pt}
\newtcolorbox{boxDef}[1][]{title={{\emoji📘} تعريف\ifx\\#1\\\else ~#1\fi :}}
\newtcolorbox{boxExercise}[1][]{title={{\emoji🧩} تمرين\ifx\\#1\\\else ~#1\fi :}}
\newtcolorbox{boxExample}[1][]{title={{\emoji📝} مثال\ifx\\#1\\\else ~#1\fi :}}
\newtcolorbox{boxNote}[1][]{title={{\emoji✨} ملاحظة\ifx\\#1\\\else ~#1\fi :}}
\newtcolorbox{boxAttention}[1][]{title={{\emoji🔔} تنبيه\ifx\\#1\\\else ~#1\fi :}}
\newtcolorbox{boxWarning}[1][]{title={{\emoji⚡} ملاحظة هامة\ifx\\#1\\\else ~#1\fi :}}
\newtcolorbox{boxSolution}[1][]{title={{\emoji✅} حل\ifx\\#1\\\else ~#1\fi :}}
\newtcolorbox{boxSymbol}[1][]{title={{\emoji🔣} رمز\ifx\\#1\\\else ~#1\fi :}}
\newtcolorbox{boxHint}[1][]{title={{\emoji💡} تلميح\ifx\\#1\\\else ~#1\fi :}}
\else
\tcbset{colback=white, colframe=black, fonttitle=\bfseries, boxrule=0.8pt}
\newtcolorbox{boxDef}[1][]{colback=blue!5!white,colframe=blue!75!black,
  title={{\emoji📘} تعريف\ifx\\#1\\\else ~#1\fi :}}
\newtcolorbox{boxExercise}[1][]{colback=cyan!5!white,colframe=cyan!70!black,
  title={{\emoji🧩} تمرين\ifx\\#1\\\else ~#1\fi :}}
\newtcolorbox{boxExample}[1][]{colback=yellow!5!white,colframe=orange!90!black,
  title={{\emoji📝} مثال\ifx\\#1\\\else ~#1\fi :}}
\newtcolorbox{boxNote}[1][]{colback=gray!10!white,colframe=black,
  title={{\emoji✨} ملاحظة\ifx\\#1\\\else ~#1\fi :}}
\newtcolorbox{boxAttention}[1][]{colback=magenta!10!white,colframe=magenta!80!black,
  title={{\emoji🔔} تنبيه\ifx\\#1\\\else ~#1\fi :}}
\newtcolorbox{boxWarning}[1][]{colback=red!5!white,colframe=red!75!black,
  title={{\emoji⚡} ملاحظة هامة\ifx\\#1\\\else ~#1\fi :}}
\newtcolorbox{boxSolution}[1][]{colback=green!5!white,colframe=green!60!black,
  title={{\emoji✅} حل\ifx\\#1\\\else ~#1\fi :}}
\newtcolorbox{boxSymbol}[1][]{colback=purple!5!white,colframe=purple!70!black,
  title={{\emoji🔣} رمز\ifx\\#1\\\else ~#1\fi :}}
\newtcolorbox{boxHint}[1][]{colback=teal!5!white,colframe=teal!60!black,
  title={{\emoji💡} تلميح\ifx\\#1\\\else ~#1\fi :}}
\fi


\ifwhiteblack
\tcbset{simplecode/.style={ colback=white, colframe=black, boxrule=1pt, arc=2pt, left=4pt,right=4pt,top=4pt,bottom=4pt}}
\else
\tcbset{simplecode/.style={ colback=gray!5, colframe=black!50, boxrule=0.4pt, arc=2pt, left=4pt,right=4pt,top=4pt,bottom=4pt}}
\fi
% Compact code boxes: tighter line spacing (the preamble uses onehalfspacing) and small vertical gaps.
% Must come AFTER the simplecode definition above.
\tcbset{simplecode/.append style={fontupper=\setstretch{0.95}, top=4pt, bottom=4pt, before skip=6pt, after skip=6pt}}
\newenvironment{boxCode}{\begin{tcolorbox}[simplecode]}{\end{tcolorbox}}

% Code cell used inside tables (keeps code left-to-right inside RTL tables).
% Use inside a fixed-width p{...} column: \begin{pycell}\begin{minted}...\end{minted}\end{pycell}
\newtcolorbox{pycell}{frame empty,colback=white,boxsep=0pt,left=2pt,right=2pt,top=2pt,bottom=2pt,width=\linewidth}

\newcolumntype{C}[1]{>{\centering\arraybackslash}p{#1}}

\makeatletter
\DTMnewdatestyle{dotted}{%
  \renewcommand{\DTMdisplaydate}[4]{\two@digits{##3}.\two@digits{##2}.##1}%
  \renewcommand{\DTMDisplaydate}{\DTMdisplaydate}%
}
\makeatother
\DTMsetdatestyle{dotted}

% redefine spaces after titles
\makeatletter
\renewcommand{\@maketitle}{%
  \begin{center}
    {\huge \bfseries \@title \par}%
    \vskip 0.2em % space between title and author
    {\large \@author \par}%
    % \vskip 0.2em % space between author and date
    % {\normalsize \@date \par}%
    {\small \hfill \begin{english}\DTMtoday\end{english} \par}%
  \end{center}
}
\makeatother

\fancyhf{} % clear default
\fancypagestyle{plain}{
  \fancyhf{}
  \fancyhead[L]{مدرسة التسامح الشاملة}
  % \fancyhead[L]{\includegraphics[height=1cm]{../../../images/logoTasamoh.png}}
  \fancyhead[R]{الأستاذ محمود اغبارية}
  \fancyfoot[C]{\thepage}
}

\fancyhead[L]{مدرسة التسامح الشاملة}
\fancyhead[R]{الأستاذ محمود اغبارية}
\fancyfoot[C]{\thepage}

\setcounter{tocdepth}{3} % only section subsection and subsubsection in TOC

\newcommand{\importpdfpage}[2]{%
    \thispagestyle{empty}
    \newgeometry{left=0cm,right=0cm,top=0cm,bottom=0cm}
    \includegraphics[page=#2,width=0.95\paperwidth,height=0.95\paperheight,keepaspectratio]{#1}
    \restoregeometry
}

\newcommand{\insertFullImg}[1]{%
    \noindent
    \makebox[\textwidth][c]{
        \includegraphics[width=0.9\paperwidth,keepaspectratio]{#1}%
    }%
}%


\ifwithsols
\title{حل ورقة مراجعة للصف العاشر 10 - الشروط}
\else
\title{ورقة مراجعة للصف العاشر 10 - الشروط}
\fi

\begin{document}

\maketitle
\thispagestyle{fancy}

%====================================================================
\section*{القسم الأول: التعابير المنطقية والشروط المركبة}
%====================================================================

\begin{enumerate}[label=\textbf{تمرين \arabic*:}, wide=0pt, leftmargin=*, itemsep=2.2em]

% ---------------------------------------------------------------- 1
\item لكل تعبير منطقي في الجدول التالي، ضعوا إشارة \textenglish{X} في العمود الملائم لنتيجته: \textenglish{True} أم \textenglish{False}.

\begin{english}
\renewcommand{\arraystretch}{1.6}
\begin{longtable}{|C{0.8cm}|p{9cm}|C{2cm}|C{2cm}|}
\hline
\textbf{\#} & \textarabic{\textbf{التعبير المنطقي}} & \textbf{True} & \textbf{False} \\
\hline
\endhead
1 & \texttt{3 < 5 and 9 > 10} & & \\ \hline
2 & \texttt{8 < 7 or 16 > 15} & & \\ \hline
3 & \texttt{4 <= 3 and 2 == 2} & & \\ \hline
4 & \texttt{10 <= 10 or 5 != 5} & & \\ \hline
5 & \texttt{12 < 10 and 5 >= 3} & & \\ \hline
6 & \texttt{9 != 9 and 6 > 5} & & \\ \hline
7 & \texttt{(3 > 5 and 9 < 10) or (6 == 6)} & & \\ \hline
8 & \texttt{(4 < 2 and 8 > 7) or (5 < 5)} & & \\ \hline
9 & \texttt{(7 > 6 and 12 < 11) and (8 > 8)} & & \\ \hline
10 & \texttt{(not(2 < 1) and 6 > 5) or not(5 != 10)} & & \\ \hline
\end{longtable}
\end{english}

\ifwithsols
\begin{boxSolution}
\begin{english}
\renewcommand{\arraystretch}{1.2}
\begin{tabular}{rcl}
1 & $\rightarrow$ & \texttt{False} \\
2 & $\rightarrow$ & \texttt{True} \\
3 & $\rightarrow$ & \texttt{False} \\
4 & $\rightarrow$ & \texttt{True} \\
5 & $\rightarrow$ & \texttt{False} \\
6 & $\rightarrow$ & \texttt{False} \\
7 & $\rightarrow$ & \texttt{True} \\
8 & $\rightarrow$ & \texttt{False} \\
9 & $\rightarrow$ & \texttt{False} \\
10 & $\rightarrow$ & \texttt{True} \\
\end{tabular}
\end{english}
\end{boxSolution}

\fi

\clearpage
% ---------------------------------------------------------------- 2
\item ضعوا دائرة حول \textbf{المخرج} الملائم لكل أمر: \textenglish{True} أم \textenglish{False}.

\begin{english}
\renewcommand{\arraystretch}{1.6}
\begin{longtable}{|C{0.8cm}|p{10cm}|C{3cm}|}
\hline
\textbf{\#} & \textarabic{\textbf{الأمر}} & \textarabic{\textbf{المخرج}} \\
\hline
\endhead
1 & \texttt{print(4 < 2 and 8 > 7)} & True / False \\ \hline
2 & \texttt{print(5 < 4 or 11 > 10)} & True / False \\ \hline
3 & \texttt{print(9 < 8 and 15 > 14)} & True / False \\ \hline
4 & \texttt{print(6 < 5 and 13 > 12)} & True / False \\ \hline
5 & \texttt{print(1 < 0 and 4 > 3)} & True / False \\ \hline
6 & \texttt{print(8 < 7 or 16 > 15)} & True / False \\ \hline
7 & \texttt{print(10 < 9 and 20 > 19)} & True / False \\ \hline
8 & \texttt{print((4 < 2 and 8 > 7) or not(5 == 5))} & True / False \\ \hline
9 & \texttt{print((not(2 < 1) and 6 > 5) or not(10 == 10))} & True / False \\ \hline
10 & \texttt{print((not(9 > 8) and not(15 < 14)) or (11 < 11))} & True / False \\ \hline
11 & \texttt{print((8 > 7 or not(16 < 15)) or (14 == 14))} & True / False \\ \hline
12 & \texttt{print((not(10 < 9) and 20 > 19) or not(15 == 15))} & True / False \\ \hline
\end{longtable}
\end{english}

\ifwithsols
\begin{boxSolution}
\begin{english}
\renewcommand{\arraystretch}{1.2}
\begin{tabular}{rcl}
1 & $\rightarrow$ & \texttt{False} \\
2 & $\rightarrow$ & \texttt{True} \\
3 & $\rightarrow$ & \texttt{False} \\
4 & $\rightarrow$ & \texttt{False} \\
5 & $\rightarrow$ & \texttt{False} \\
6 & $\rightarrow$ & \texttt{True} \\
7 & $\rightarrow$ & \texttt{False} \\
8 & $\rightarrow$ & \texttt{False} \\
9 & $\rightarrow$ & \texttt{True} \\
10 & $\rightarrow$ & \texttt{False} \\
11 & $\rightarrow$ & \texttt{True} \\
12 & $\rightarrow$ & \texttt{True} \\
\end{tabular}
\end{english}
\end{boxSolution}

\fi

\end{enumerate}

\clearpage
%====================================================================
\section*{القسم الثاني: كتابة برامج بشروط مركبة (اختياري)}
%====================================================================

\begin{enumerate}[resume, label=\textbf{تمرين \arabic*:}, wide=0pt, leftmargin=*, itemsep=2.2em]

% ---------------------------------------------------------------- 3
\item \textbf{شراء حاسوب نقّال:} اكتبوا برنامجًا \textbf{يستقبل} المبلغ المالي الذي تم توفيره وعدد أشهر التوفير. إذا كان المبلغ الذي تم توفيره \textbf{أكثر من} 4000 ش.ج \textbf{أو} عدد الأشهر \textbf{أكثر من} 10، \textbf{يطبع} البرنامج عبارة توضيحية أنّ بإمكانه شراء حاسوب نقّال.

\ifwithsols
\begin{boxSolution}
\begin{english}
\begin{minted}{python}
saved_amount = float(input("Enter the amount of money saved: "))
months = int(input("Enter the number of saving months: "))

if saved_amount > 4000 or months > 10:
    print("You can buy a laptop!")
\end{minted}
\end{english}
\end{boxSolution}

\filbreak
\fi

% ---------------------------------------------------------------- 4
\item \textbf{مجلّد الصور:} اكتبوا برنامجًا \textbf{يستقبل} مساحة الذاكرة المتاحة في مجلّد الصور المحوسب. إذا كانت القيمة المُدخلة \textbf{أقل من} 2 ميغابايت (\textenglish{MB})، \textbf{يطبع} البرنامج عبارة توضيحية أنّ المكان ممتلئ وأنّه يجب حذف صور.

\ifwithsols
\begin{boxSolution}
\begin{english}
\begin{minted}{python}
free_space = float(input("Enter the free space in the picture folder (in megabytes): "))

if free_space < 2:
    print("Out of space! Delete some pictures.")
\end{minted}
\end{english}
\end{boxSolution}

\filbreak
\fi

% ---------------------------------------------------------------- 5
\item \textbf{قطار الجبال:} اكتبوا برنامجًا \textbf{يستقبل} طول شخص (بالسنتيمتر) وعمره. إذا كان طول الشخص 130 سم \textbf{وما فوق} \textbf{و}عمره 13 سنة \textbf{وما فوق}، \textbf{يطبع} البرنامج عبارة توضيحية أنّ بإمكانه الدخول إلى قطار الجبال.

\ifwithsols
\begin{boxSolution}
\begin{english}
\begin{minted}{python}
height = int(input("Enter the person's height: "))
age = int(input("Enter the person's age: "))

if height >= 130 and age >= 13:
    print("You can go on the roller coaster!")
\end{minted}
\end{english}
\end{boxSolution}

\filbreak
\fi

% ---------------------------------------------------------------- 6
\item \textbf{رخصة السياقة:} يحصل الطالب على رخصة سياقة إذا تحققت الشروط \textbf{الثلاثة} التالية:

\begin{itemize}[itemsep=0.3em]
\item عمره \textbf{أكبر من} 17.
\item نجح في امتحان السياقة \textbf{النظري} (\textenglish{theory}).
\item نجح في امتحان السياقة \textbf{العملي}.
\end{itemize}

اكتبوا برنامجًا \textbf{يستقبل} المعطيات الثلاثة. إذا كان الطالب لا يستحق رخصة سياقة، \textbf{يطبع} البرنامج السبب، وإلا \textbf{يطبع} عبارة تهنئة ملائمة.

\ifwithsols
\begin{boxSolution}
\begin{english}
\begin{minted}{python}
age = int(input("Enter the age of the applicant: "))
has_passed_theory_test = input("Has the applicant passed the theory test? (yes/no): ").lower()
has_passed_driving_test = input("Has the applicant passed the driving test? (yes/no): ").lower()

if age <= 17:
    print("Applicant is not eligible for a driving license due to age.")
elif has_passed_theory_test != "yes":
    print("Applicant is not eligible for a driving license as they have not passed the theory test.")
elif has_passed_driving_test != "yes":
    print("Applicant is not eligible for a driving license as they have not passed the driving test.")
else:
    print("Applicant is eligible for a driving license. Congratulations!")
\end{minted}
\end{english}
\end{boxSolution}

\begin{boxNote}
العملية \textenglish{lower()} تحوّل جميع أحرف السلسلة النصية إلى أحرف صغيرة. سنتحدث عن ذلك بشكل موسّع فيما بعد.
\end{boxNote}

\fi

\end{enumerate}

\clearpage
%====================================================================
\section*{القسم الثالث: الشرط المتداخل}
%====================================================================

\begin{enumerate}[resume, label=\textbf{تمرين \arabic*:}, wide=0pt, leftmargin=*, itemsep=2.2em]

% ---------------------------------------------------------------- 7
\item \textbf{منزلة واحدة:} اكتبوا برنامجًا \textbf{يستقبل} عددًا صحيحًا. يفحص البرنامج إذا كان العدد \textbf{موجبًا} ومركّبًا من \textbf{منزلة واحدة فقط}، و\textbf{يطبع} عبارة توضيحية ملائمة. \textbf{يجب استخدام الشرط المتداخل.}

\textbf{تحدٍّ:} غيّروا البرنامج بحيث يفحص كل عدد صحيح (\textbf{موجبًا} أو \textbf{سالبًا}).

\begin{boxAttention}[- نقطة للتفكير]
هل توجد طريقة أخرى لحل التمرين؟
\end{boxAttention}

\ifwithsols
\begin{boxSolution}[1]
\begin{english}
\begin{minted}{python}
number = int(input("Enter a number: "))

if number > 0:
    if number <= 9:
        print(number, "is a single digit number.")
    else:
        print(number, "is not a single digit number.")
else:
    print(number, "is not a positive number.")
\end{minted}
\end{english}
\end{boxSolution}

\begin{boxSolution}[بديل: شرط مركب]
\begin{english}
\begin{minted}{python}
number = int(input("Enter a number: "))

if number > 0 and number <= 9:
    print(number, "is a single digit number.")
else:
    print(number, "is not a single digit number.")
\end{minted}
\end{english}
\end{boxSolution}

\filbreak

\begin{boxSolution}[التحدي]
\begin{english}
\begin{minted}{python}
number = int(input("Enter a number: "))

if number >= -9:
    if number <= 9:
        print(number, "is a single digit number.")
    else:
        print(number, "is not a single digit number.")
else:
    print(number, "is not a single digit number.")
\end{minted}
\end{english}
\end{boxSolution}

\fi
\filbreak

% ---------------------------------------------------------------- 8
\item \textbf{العدد الأكبر:} اكتبوا برنامجًا \textbf{يستقبل} ثلاثة أعداد. يفحص البرنامج أيّها \textbf{الأكبر} و\textbf{يطبعه}. \textbf{يجب استخدام الشرط المتداخل.}

\ifwithsols
\begin{boxSolution}[1]
\begin{english}
\begin{minted}{python}
# Read three numbers from the user
number1 = float(input("Enter first number: "))
number2 = float(input("Enter second number: "))
number3 = float(input("Enter third number: "))

# Find the largest number using nested conditions
if number1 >= number2:
    if number1 >= number3:
        max_num = number1
    else:
        max_num = number3
else:
    if number2 >= number3:
        max_num = number2
    else:
        max_num = number3

print("The largest number is:", max_num)
\end{minted}
\end{english}
\end{boxSolution}

\filbreak

\begin{boxSolution}[بديل: شروط مركبة ومتعددة]
\begin{english}
\begin{minted}{python}
number1 = float(input("Enter first number: "))
number2 = float(input("Enter second number: "))
number3 = float(input("Enter third number: "))

if number1 >= number2 and number1 >= number3:
    print(number1, "is the largest number.")
elif number2 >= number1 and number2 >= number3:
    print(number2, "is the largest number.")
else:
    print(number3, "is the largest number.")
\end{minted}
\end{english}
\end{boxSolution}

\fi
\filbreak

% ---------------------------------------------------------------- 9
\item \textbf{تعليمات عبور الشارع:} قبل أن نعبر الشارع بجانب المدرسة يجب أن يتحقق الشرطان التاليان:

\begin{itemize}[itemsep=0.3em]
\item على الإشارة الضوئية أن تكون \textbf{خضراء}.
\item يجب أن يتواجد \textbf{على الأقل طالبان} يريدان عبور الشارع.
\end{itemize}

وُضعت إشارة محوسبة على الرصيف تعرض إذا كان مسموحًا عبور الشارع. اكتبوا عملية \textbf{تستقبل} من المستخدم المعلومات المطلوبة و\textbf{تطبع} على الإشارة عبارة توضيحية ملائمة. \textbf{يجب استخدام الشرط المتداخل.}

استدعوا العملية عدة مرات بمُدخلات مختلفة، بحيث تظهر كل عبارة توضيحية ممكنة.

\ifwithsols
\begin{boxSolution}
\begin{english}
\begin{minted}{python}
def crossing_sign():
    is_green_light = input("Is the traffic light green? (yes/no): ")

    if is_green_light == "yes":
        num_students = int(input("Enter the number of students on the sidewalk: "))
        if num_students >= 2:
            print("You can cross the road")
        else:
            print("You cannot cross the road. Wait for another friend")
    else:
        print("You cannot cross the road. No green light")

# Call the function
crossing_sign()
\end{minted}
\end{english}
\end{boxSolution}

\fi
\filbreak

% ---------------------------------------------------------------- 10
\item \textbf{شروط القبول لتخصص علوم الحاسوب:} اكتبوا عملية \textbf{تستقبل} من المستخدم علامات طالب في ثلاثة مواضيع، و\textbf{تطبع} عبارة توضيحية إذا تم قبول الطالب أم لا، بحسب شروط القبول التالية:

\begin{itemize}[itemsep=0.3em]
\item اللغة الإنجليزية: علامة 85 \textbf{فما فوق}.
\item الرياضيات: علامة 80 \textbf{فما فوق}.
\item العلوم: علامة 75 \textbf{فما فوق}.
\end{itemize}

(يُقبل الطالب فقط إذا تحققت الشروط \textbf{الثلاثة}.) \textbf{يجب استخدام الشرط المتداخل.}

\ifwithsols
\begin{boxSolution}
\begin{english}
\begin{minted}{python}
def admission():
    english_grade = int(input("Enter grade in English: "))
    math_grade = int(input("Enter grade in math: "))
    sciences_grade = int(input("Enter grade in science: "))

    if english_grade >= 85:
        if math_grade >= 80:
            if sciences_grade >= 75:
                print("Welcome! The student is accepted")
            else:
                print("We apologize, the student is not accepted.")
        else:
            print("We apologize, the student is not accepted.")
    else:
        print("We apologize, the student is not accepted.")

# Call the function
admission()
\end{minted}
\end{english}
\end{boxSolution}

\fi
\filbreak

% ---------------------------------------------------------------- 11
\item \textbf{من مضاعفات العدد 6:} اكتبوا عملية \textbf{تستقبل} من المستخدم عددًا صحيحًا \textbf{موجبًا}. تفحص العملية إذا كان العدد من مضاعفات العدد 6، ثم \textbf{تطبع} عبارة توضيحية ملائمة.

\ifwithsols
\begin{boxSolution}[1]
\begin{english}
\begin{minted}{python}
def multiple_of_six():
    num = int(input("Enter a number: "))

    if num % 2 == 0:
        if num % 3 == 0:
            print(num, "is a multiple of 6")
        else:
            print(num, "is not a multiple of 6")
    else:
        print(num, "is not a multiple of 6")

# Call the function
multiple_of_six()
\end{minted}
\end{english}
\end{boxSolution}

\begin{boxSolution}[بديل]
\begin{english}
\begin{minted}{python}
def multiple_of_six():
    num = int(input("Enter a number: "))

    if num % 6 == 0:
        print(num, "is a multiple of 6")
    else:
        print(num, "is not a multiple of 6")

# Call the function
multiple_of_six()
\end{minted}
\end{english}
\end{boxSolution}

\fi
\filbreak

% ---------------------------------------------------------------- 12
\item \textbf{تحويل مقطع برنامج:} أمامكم مقطع برنامج فيه شرط \textbf{مركّب}:

\begin{boxCode}
\begin{english}
\begin{minted}{python}
if num <= 10 and num > 4:
    print("Correct")
\end{minted}
\end{english}
\end{boxCode}

اكتبوا مقطع برنامج يؤدي نفس الغرض ولكن بواسطة \textbf{الشرط المتداخل}.

\ifwithsols
\begin{boxSolution}
\begin{english}
\begin{minted}{python}
if num <= 10:
    if num > 4:
        print("Correct")
\end{minted}
\end{english}
\end{boxSolution}

\fi
\filbreak

% ---------------------------------------------------------------- 13
\item \textbf{تحويل عملية:} أمامكم عملية فيها شرط \textbf{مركّب}:

\begin{boxCode}
\begin{english}
\begin{minted}{python}
def ex7():
    a = int(input("Enter first number: "))
    b = int(input("Enter second number: "))
    if a > b and a < 20:
        print("True")
    else:
        print("False")

ex7()
\end{minted}
\end{english}
\end{boxCode}

اكتبوا العملية من جديد بحيث تؤدي نفس الغرض ولكن بواسطة \textbf{الشرط المتداخل}.

\ifwithsols
\begin{boxSolution}
\begin{english}
\begin{minted}{python}
def ex7():
    a = int(input("Enter first number: "))
    b = int(input("Enter second number: "))
    if a > b:
        if a < 20:
            print("True")
        else:
            print("False")
    else:
        print("False")

ex7()
\end{minted}
\end{english}
\end{boxSolution}

\fi
\filbreak

% ---------------------------------------------------------------- 14
\item \textbf{برنامج لأحجية} (تمرين تحدٍّ، يُنفَّذ في أزواج): عليكم أن تكتبوا برنامجًا لأحجية معيّنة.

\begin{enumerate}[label=\arabic*., leftmargin=*, itemsep=0.5em]
\item اختاروا موضوعًا معيّنًا واكتبوا 5--10 أسئلة عنه.
\item اكتبوا برنامجًا يطرح الأسئلة الواحد تلو الآخر على المشترك في الأحجية و\textbf{يستقبل} إجابته:
\begin{itemize}[itemsep=0.3em]
\item إذا كانت إجابته صحيحة، \textbf{يطبع} البرنامج عبارة توضيحية ملائمة.
\item يحصل المشترك على \textbf{20 نقطة} عن كل إجابة صحيحة.
\item في نهاية الأحجية، \textbf{يطبع} البرنامج عدد النقاط التي جمعها المشترك.
\end{itemize}
\end{enumerate}

\ifwithsols
\begin{boxSolution}[نموذج لحل (3 أسئلة، يمكن توسيعه)]
\begin{english}
\begin{minted}{python}
score = 0

answer = input("Which keyword defines a function in Python? ")
if answer == "def":
    print("Correct!")
    score = score + 20
else:
    print("Wrong answer.")

answer = input("Which operator checks if two values are equal? ")
if answer == "==":
    print("Correct!")
    score = score + 20
else:
    print("Wrong answer.")

answer = input("Which keyword is short for 'else if' in Python? ")
if answer == "elif":
    print("Correct!")
    score = score + 20
else:
    print("Wrong answer.")

print("Your final score is:", score)
\end{minted}
\end{english}
\end{boxSolution}
\fi

\end{enumerate}

\vspace{1cm}
\begin{flushleft}
أرجو لكم وقتًا ممتعًا.

الأستاذ محمود اغبارية.
\end{flushleft}

\end{document}
